\documentclass[11pt]{amsart}
\usepackage[margin=1.1in]{geometry}
\usepackage{amssymb,amsmath,amsthm}
\usepackage{booktabs}
\usepackage{hyperref}

\newtheorem{theorem}{Theorem}[section]
\newtheorem{lemma}[theorem]{Lemma}
\newtheorem{proposition}[theorem]{Proposition}
\newtheorem{corollary}[theorem]{Corollary}
\theoremstyle{definition}
\newtheorem{definition}[theorem]{Definition}
\newtheorem{remark}[theorem]{Remark}
\newtheorem*{mainthm}{Theorem}

\newcommand{\F}{\mathbb{F}_3}
\newcommand{\Z}{\mathbb{Z}}
\newcommand{\wt}{\operatorname{wt}}
\newcommand{\Index}{\operatorname{Index}}
\newcommand{\cF}{\mathcal{F}}
\newcommand{\isq}{\sqrt{-3}}

\title[The index of a self-dual ternary code]{The index of a self-dual ternary code is divisible by \texorpdfstring{$24$}{24}: four proofs}
\author{[Draft --- authors to be added]}
\date{Draft, 26 September 2026}

\begin{document}
\begin{abstract}
Gaiotto and Johnson-Freyd conjectured that the index of every self-dual ternary code $C$, the
number of even minus the number of odd maximal codewords, is divisible by $24$. They proved this for lengths divisible by $24$
using Borcherds' theorem on theta series of even unimodular lattices, and remarked that they did not
know an elementary proof. We record several proofs for all lengths. The main one is finite and
elementary: a single Poisson summation over $C$, with two marked coordinates, shows that the
maximal codewords, counted with signs, form a $2$-design: $\sum_{\varepsilon}s(\varepsilon)\,\varepsilon_i\varepsilon_j=0$
for $i\ne j$. Both factors $8$ and $3$ of $24$ follow from this by short counting arguments.
We also give a proof of the design identity by modular forms, a proof of the factor $8$ by Kneser
neighbours, and a neighbour version of the Poisson argument. The same computation with four marked
coordinates gives an identity that explains why the design property stops at strength two.
\end{abstract}
\maketitle

\section{Introduction}

A ternary code of length $n$ is a subspace $C\subseteq\F^n$. It is \emph{self-dual} if
$C=C^\perp$ for the standard form $x\cdot y=\sum_k x_ky_k$. A codeword is \emph{maximal} if none of
its entries is $0$, and it is \emph{even} or \emph{odd} according to the parity of the number of
its entries equal to $1$. The \emph{index} of a self-dual code is
\[
  \Index(C)=\#\{\text{even maximal codewords}\}-\#\{\text{odd maximal codewords}\}.
\]
Gaiotto and Johnson-Freyd \cite{GJF} showed that $\Index(C)$ is the supersymmetric (Witten) index
of a holomorphic $\mathcal N=1$ SVOA of central charge $n$ attached to $C$, and made the following
conjecture.

\begin{quote}
\textbf{Conjecture} (\cite[Conj.~1.4]{GJF}; numbering of the arXiv version). The index of any
self-dual ternary code is divisible by $24$, even if the length is merely divisible by $12$. For all
lengths $12k$ there exists a self-dual ternary code of index $24$.
\end{quote}

They proved the first sentence when $24\mid n$ \cite[Thm.~1.3]{GJF}, and deduced $12\mid\Index(C)$
when $12\mid n$. About their proof they write: ``Our proof relies on some nontrivial facts about Theta
functions of lattices, and we do not know if there is an elementary proof.'' This note records
proofs of the divisibility statement in all lengths.

\begin{mainthm}
For every self-dual ternary code $C$ of positive length, $24\mid\Index(C)$.
\end{mainthm}

The hypothesis $n\ge1$ is needed: for $n=0$ the empty word is an even maximal codeword, so
$\Index=1$. It was made explicit by the formalization in Lean (Remark~\ref{rem:lean}).

The routes are as follows.
\begin{itemize}
\item[(0)] \emph{GJF, $24\mid n$} (\S\ref{sec:gjf}): the two even neighbours of the lattice $\Lambda(C)$,
  and Borcherds' theorem that the $\Delta^{k}$-coefficient of the theta series of an even unimodular
  lattice of rank $24k$ is divisible by $24$.
\item[(A)] \emph{Kneser neighbours} (\S\ref{sec:A}): a Fourier identity on $\F^n$ shows
  $\Index(C)+\Index(C')\equiv0\pmod8$ for neighbouring codes, and the neighbour graph is connected.
  This gives $8\mid\Index$. The factor $3$ still comes from route (0).
\item[(B)] \emph{Signed $2$-design, modular} (\S\ref{sec:B}): the design identity
  (Theorem~\ref{thm:design}) from a theta series with harmonic coefficients and
  $M_2(\mathrm{SL}_2(\Z))=0$. Proposition~\ref{prop:design-suffices} then gives $8$ and $3$.
\item[(C)] \emph{Signed $2$-design, finite} (\S\ref{sec:C}): the design identity from one Poisson
  summation over $C$ with two marked coordinates, and again Proposition~\ref{prop:design-suffices}.
  This proof is elementary and self-contained. A neighbour version (C$'$) is in
  Remark~\ref{rem:Cprime}.
\end{itemize}
Routes (B) and (C) are each a complete proof of the Theorem, and so is (A) together with (0). A
comparison is in \S\ref{sec:compare}.

\section{Preliminaries}\label{sec:prelim}

Throughout, $n\ge1$.

\begin{lemma}\label{lem:basic}
Let $C\subseteq\F^n$ be self-dual.
\begin{enumerate}
\item $4\mid n$.
\item Every $u$ in a self-orthogonal code $D\subseteq D^\perp$ has $\wt(u)\equiv0\pmod3$.
\item If $C$ has a maximal codeword, then $12\mid n$.
\end{enumerate}
\end{lemma}

\begin{proof}
(1) $C$ is a totally isotropic subspace of half dimension of the standard quadratic space $\F^n$,
so $\F^n$ is an orthogonal sum of $n/2$ hyperbolic planes. Comparing discriminants in
$\F^\times/(\F^\times)^2=\{\pm1\}$ gives $1=(-1)^{n/2}$, so $n/2$ is even.
(2) $u\cdot u=\sum_ku_k^2=\wt(u)$ in $\F$, since $a^2=1$ for $a\ne0$.
(3) A maximal codeword has weight $n$, so $3\mid n$ by (2), and $4\mid n$ by (1).
\end{proof}

\subsection*{Signs and lifts}
Let $\varphi\colon\F\to\{0,\pm1\}$ be $\varphi(0)=0$, $\varphi(1)=-1$, $\varphi(2)=1$, and
$f(x)=\prod_k\varphi(x_k)$. Then $f(x)=(-1)^{\#\{k:x_k=1\}}$ for maximal $x$ and $f(x)=0$ otherwise,
so $\Index(C)=\sum_{x\in C}f(x)$.

For a maximal $x$ let $\varepsilon(x)\in\{\pm1\}^n$ be its lift, $\varepsilon_k=-\varphi(x_k)$; so
$\varepsilon_k=1$ iff $x_k=1$, and $\varepsilon(x)\equiv x\pmod 3$ coordinatewise. Put
$s(\varepsilon)=(-1)^{\#\{k:\varepsilon_k=1\}}=f(x)$, and let $\cF=\cF(C)\subset\{\pm1\}^n$ be the
set of lifts of the maximal codewords. Then
\begin{itemize}
\item $\Index(C)=\sum_{\varepsilon\in\cF}s(\varepsilon)$;
\item $\cF=-\cF$ and $s(-\varepsilon)=s(\varepsilon)$, since $n$ is even;
\item $\varepsilon\cdot\varepsilon'\equiv0\pmod3$ for $\varepsilon,\varepsilon'\in\cF$ (integer dot
product), because it reduces to the $\F$-product of two codewords.
\end{itemize}

\subsection*{Fourier analysis on \texorpdfstring{$\F^n$}{F\_3\textasciicircum n}}
Let $\omega=e^{2\pi i/3}$ and $\isq:=\omega-\omega^2=i\sqrt3$. For $h\colon\F^n\to\mathbb C$ put
$\hat h(u)=\sum_xh(x)\,\omega^{u\cdot x}$.

\begin{lemma}\label{lem:fourier}
\begin{enumerate}
\item \emph{(Poisson.)} For every linear code $D\subseteq\F^n$,
  $\displaystyle\sum_{x\in D^\perp}h(x)=|D|^{-1}\sum_{u\in D}\hat h(u)$.
\item If $h(x)=\prod_kh_k(x_k)$, then $\hat h(u)=\prod_k\hat h_k(u_k)$.
\item $\hat\varphi=\isq\cdot\varphi$, and the indicator $\iota=[\,\cdot\ne0\,]$ has
  $\hat\iota(b)=3[b=0]-1$. Consequently $\hat f=(\isq)^nf=3^{n/2}f$ when $4\mid n$.
\end{enumerate}
\end{lemma}

\begin{proof}
(1) $\sum_{u\in D}\hat h(u)=\sum_xh(x)\sum_{u\in D}\omega^{u\cdot x}=|D|\sum_{x\in D^\perp}h(x)$.
(2) is clear. (3) $\hat\varphi(b)=-\omega^b+\omega^{2b}$ is $0,\ \omega^2-\omega=-\isq,\ \omega-\omega^2=\isq$
for $b=0,1,2$. And $\hat\iota(b)=\omega^b+\omega^{2b}$ is $2$ for $b=0$ and $-1$ otherwise.
Finally $(\isq)^2=-3$.
\end{proof}

\section{The signed \texorpdfstring{$2$}{2}-design and why it suffices}\label{sec:design}

\begin{theorem}[signed $2$-design]\label{thm:design}
For every self-dual ternary code $C$ and all coordinates $i,j$,
\[
  \sum_{\varepsilon\in\cF(C)}s(\varepsilon)\,\varepsilon_i\varepsilon_j=\Index(C)\,\delta_{ij}.
\]
Equivalently, $\sum_{\varepsilon\in\cF}s(\varepsilon)(\varepsilon\cdot\mu)^2=\Index(C)\,|\mu|^2$ for
all $\mu\in\mathbb R^n$.
\end{theorem}

The diagonal case $i=j$ is trivial, since $\varepsilon_i^2=1$. Proofs of the off-diagonal case are
in \S\ref{sec:C} and \S\ref{sec:B}.

\begin{proposition}\label{prop:design-suffices}
Theorem~\ref{thm:design} implies $8\mid\Index(C)$ and $3\mid\Index(C)$.
\end{proposition}

\begin{proof}
If $C$ has no maximal codeword, $\Index(C)=0$. Otherwise $12\mid n$ by Lemma~\ref{lem:basic}.

\emph{The factor $8$.} Fix distinct coordinates $a,b,c$. For $\varepsilon\in\{\pm1\}^n$,
$[\varepsilon_a=\varepsilon_b=\varepsilon_c]=\tfrac14(1+\varepsilon_a\varepsilon_b+\varepsilon_a\varepsilon_c
+\varepsilon_b\varepsilon_c)$. Multiplying by $s(\varepsilon)$, summing over $\cF$ and using
Theorem~\ref{thm:design},
\[
  \sum_{\varepsilon\in\cF,\ \varepsilon_a=\varepsilon_b=\varepsilon_c}s(\varepsilon)=\tfrac14\Index(C).
\]
The set summed over is stable under $\varepsilon\mapsto-\varepsilon$, which has no fixed points
and preserves $s$. So the left side is even, and $8\mid\Index(C)$.

\emph{The factor $3$.} Fix $\varepsilon^0\in\cF$. For every $x\in\Z^n$ with $x\equiv\varepsilon^0
\pmod3$ and every $\varepsilon\in\cF$ we have $\varepsilon\cdot x\equiv\varepsilon\cdot\varepsilon^0
\equiv0\pmod3$, so $\sum_{\cF}s(\varepsilon)(\varepsilon\cdot x)^2\in9\Z$. By
Theorem~\ref{thm:design} this sum is $\Index(C)\,|x|^2$. Take $x=\varepsilon^0$, with $|x|^2=n$, and
$x'=\varepsilon^0+3\varepsilon^0_1e_1$, with $|x'|^2=n-1+16=n+15$. Then $9$ divides both
$\Index(C)\,n$ and $\Index(C)(n+15)$, hence $9\mid15\,\Index(C)$, and $3\mid\Index(C)$.
\end{proof}

\begin{remark}\label{rem:golay}
For the extended ternary Golay code the $24$ maximal codewords are all even, and
Theorem~\ref{thm:design} says that the $12$ of them taken up to sign are the rows of a Hadamard
matrix of order $12$. This is the classical relation between the Golay code and the Hadamard matrix
of order $12$; see e.g.\ \cite{AHM} for Hadamard matrices formed by codewords of ternary self-dual
codes. Theorem~\ref{thm:design} is a signed version of this property that holds for every
self-dual ternary code.
\end{remark}

\section{Route C: a finite proof of the design identity}\label{sec:C}

Fix distinct coordinates $i,j$ and define
\[
  g_{ij}(x)=\prod_{k\ne i,j}\varphi(x_k)\cdot[x_i\ne0]\,[x_j\ne0],\qquad G_{ij}(X)=\sum_{x\in X}g_{ij}(x).
\]
For a maximal $x$ with lift $\varepsilon$ we have $s(\varepsilon)\varepsilon_i\varepsilon_j
=f(x)\varphi(x_i)\varphi(x_j)=g_{ij}(x)$, since $\varphi(x_i)^2=\varphi(x_j)^2=1$. And $g_{ij}$
vanishes off the maximal vectors. Hence
\[
  G_{ij}(C)=\sum_{\varepsilon\in\cF(C)}s(\varepsilon)\,\varepsilon_i\varepsilon_j .
\]

\begin{lemma}\label{lem:ghat}
$\displaystyle\hat g_{ij}(u)=(\isq)^{\,n-2}\prod_{k\ne i,j}\varphi(u_k)\cdot\bigl(3[u_i=0]-1\bigr)\bigl(3[u_j=0]-1\bigr).$
\end{lemma}

\begin{proof}
Lemma~\ref{lem:fourier}(2) and (3), with $h_k=\varphi$ for $k\ne i,j$ and $h_i=h_j=\iota$.
\end{proof}

\begin{lemma}[pointwise identity]\label{lem:pointwise}
Suppose $3\mid n$ and $D\subseteq\F^n$ is self-orthogonal. Then $\hat g_{ij}(u)=(\isq)^{n-2}g_{ij}(u)$
for every $u\in D$.
\end{lemma}

\begin{proof}
If $u_k=0$ for some $k\ne i,j$, both sides vanish. Otherwise $\wt(u)=n-2+[u_i\ne0]+[u_j\ne0]$. By
Lemma~\ref{lem:basic}(2), $\wt(u)\equiv0\pmod3$; as $3\mid n$, the weights $n-1$ and $n-2$ are
excluded, so $u_i\ne0$ and $u_j\ne0$. Then the last two factors in Lemma~\ref{lem:ghat} are
$(-1)(-1)=1$, and $\prod_{k\ne i,j}\varphi(u_k)=g_{ij}(u)$.
\end{proof}

\begin{theorem}\label{thm:C}
For every self-dual ternary code $C$ and $i\ne j$, $G_{ij}(C)=0$. Hence Theorem~\ref{thm:design}
holds, and by Proposition~\ref{prop:design-suffices}, $24\mid\Index(C)$.
\end{theorem}

\begin{proof}
If $3\nmid n$, $C$ has no maximal codewords and $G_{ij}(C)=0$. Let $3\mid n$. Poisson summation
(Lemma~\ref{lem:fourier}(1)) with $D=C=C^\perp$, $|C|=3^{n/2}$, and Lemma~\ref{lem:pointwise} give
\[
  G_{ij}(C)=3^{-n/2}\sum_{u\in C}\hat g_{ij}(u)=3^{-n/2}(\isq)^{n-2}\,G_{ij}(C)=\pm\tfrac13\,G_{ij}(C).
\]
Hence $G_{ij}(C)=0$.
\end{proof}

\begin{remark}[what the proof uses]\label{rem:uses}
Only three facts enter: Poisson summation over $C=C^\perp$; the one-variable transforms of
$\varphi$ and $\iota$; and that the weights of self-orthogonal codes are divisible by $3$. The sign
of $(\isq)^{n-2}$ is irrelevant, so the proof does not even use $4\mid n$ (for $4\mid n$ the factor is
$-\frac13$). In contrast, $f$ itself is a Fourier eigenfunction with eigenvalue $3^{n/2}=|C|$, so
Poisson summation over $C$ only returns $\Index(C)=\Index(C)$. Replacing $\varphi$ by $\iota$ at two
coordinates multiplies the eigenvalue by $-\frac13$ on every word a self-orthogonal code can
contain, and $0$ is the only fixed point of $G\mapsto-G/3$.
\end{remark}

\begin{remark}[route C$'$: the neighbour form]\label{rem:Cprime}
Let $D\subseteq C$ have dimension $n/2-1$, so $D$ is self-orthogonal and $|D|=3^{n/2-1}$.
Lemma~\ref{lem:pointwise} and Poisson summation over $D$ give $G_{ij}(D^\perp)=-G_{ij}(D)$ when
$4\mid n$ and $3\mid n$. If $C'$ is the other self-dual code containing $D$, the coset decomposition in
the proof of Lemma~\ref{lem:A-neighbour} (which only uses that $g_{ij}$ is supported on maximal,
hence isotropic, vectors) gives $G_{ij}(D^\perp)=G_{ij}(C)+G_{ij}(C')-G_{ij}(D)$. Hence
$G_{ij}(C)+G_{ij}(C')=0$. Connectivity of the neighbour graph (Lemma~\ref{lem:connected}) and the
base code $T^{n/4}$, which has no maximal codewords, give a second proof of Theorem~\ref{thm:C}.
Unlike $G_{ij}(C)$, the value $G_{ij}(D)$ is usually nonzero, so this identity is not vacuous.
\end{remark}

\begin{remark}[four marked coordinates]\label{rem:four}
Let $K=\{a,b,c,d\}$ be four distinct coordinates, and $g_K$ the analogue of $g_{ij}$ with $\iota$ at
the four points of $K$. Then $\sum_{x\in C}g_K(x)=\sum_{\cF}s(\varepsilon)\varepsilon_a\varepsilon_b
\varepsilon_c\varepsilon_d=:Q_K(C)$. In the analogue of Lemma~\ref{lem:pointwise} a codeword that is
nonzero off $K$ can now have $0$ or $3$ zeros in $K$ (weights $n$ or $n-3$), and a word with three
zeros in $K$ contributes the factor $2^3\cdot(-1)=-8$. With $(\isq)^{n-4}=3^{n/2-2}$ for $4\mid n$,
Poisson summation gives $Q_K=\frac19(Q_K-8W_K)$, that is,
\[
  Q_K(C)=-W_K(C),\qquad W_K(C)=\sum_{\substack{u\in C,\ u_k\ne0\ (k\notin K)\\ \#\{k\in K:\,u_k=0\}=3}}
  \ \prod_{k\notin K}\varphi(u_k).
\]
So the weight-$(n-3)$ codewords whose zeros lie in $K$ are exactly the obstruction to a signed
$4$-design. For the Golay code $Q_K=\pm8\ne0$ for all $495$ sets $K$, so the signed maximal words
are in general not a $4$-design, in line with $M_4(\mathrm{SL}_2(\Z))\ne0$ in route B.
\end{remark}

\begin{remark}[relation to the literature]\label{rem:lit}
Theorem~\ref{thm:C} is an evaluation of a MacWilliams-type identity for the complete weight
enumerator of $C$ split at the two coordinates $i,j$, i.e.\ for a complete version of Ozeki's Jacobi
polynomials \cite{Oz,CM}. We have not found the design identity, or its consequences for the index,
in the literature; this search was not systematic. The observation that invariant theory of the
unsplit complete weight enumerator alone only gives $\Index\in6\Z$ or $2\Z$ (project notes, \S3(e))
is consistent with the need for two marked coordinates.
\end{remark}

\section{Route A: Kneser neighbours and the factor \texorpdfstring{$8$}{8}}\label{sec:A}

Throughout, $12\mid n$ (otherwise all indices vanish). For $X\subseteq\F^n$ write
$F(X)=\sum_{x\in X}f(x)$, so $F(C)=\Index(C)$.

\begin{lemma}\label{lem:A-poisson}
If $D\subseteq\F^n$ is linear of dimension $n/2-1$, then $F(D^\perp)=3F(D)$.
\end{lemma}

\begin{proof}
Poisson summation and $\hat f=3^{n/2}f$ (Lemma~\ref{lem:fourier}(3), using $4\mid n$) give
$F(D^\perp)=3^{-(n/2-1)}\cdot3^{n/2}F(D)$.
\end{proof}

\begin{lemma}[neighbour identity]\label{lem:A-neighbour}
Let $C\ne C'$ be self-dual with $D=C\cap C'$ of dimension $n/2-1$. Then
$\Index(C)+\Index(C')=4F(D)$.
\end{lemma}

\begin{proof}
Since $D\subseteq C=C^\perp$ and $D\subseteq C'$, we have $C+C'\subseteq D^\perp$. Both have
dimension $n/2+1$, so $D^\perp=C+C'$. Pick $c\in C\setminus D$ and $c'\in C'\setminus D$. Then
$c\cdot c'\ne0$: otherwise $c'\perp D+\langle c\rangle=C$, so $c'\in C$. The nine cosets
$D+ac+bc'$ ($a,b\in\F$) partition $D^\perp$. For $ab\ne0$ every vector $x$ of the coset has
$x\cdot x=2ab\,(c\cdot c')\ne0$, while maximal vectors have $x\cdot x=n=0$ in $\F$. So these four
cosets contribute nothing, and $F(D^\perp)=F(C)+F(C')-F(D)$. Lemma~\ref{lem:A-poisson} finishes the
proof.
\end{proof}

\begin{lemma}\label{lem:A-even}
$F(D)$ is even for every linear $D$.
\end{lemma}

\begin{proof}
$f(0)=0$, and $x\mapsto-x$ is a fixed-point-free involution of $D\setminus\{0\}$ with
$f(-x)=(-1)^nf(x)=f(x)$.
\end{proof}

\begin{lemma}[connectivity]\label{lem:connected}
Any two self-dual codes of length $n$ are joined by a chain of neighbours.
\end{lemma}

\begin{proof}
Let $C_1\ne C_2$ and pick $x\in C_2\setminus C_1$. Since $x\notin C_1=C_1^\perp$, the code
$C_1\cap x^\perp$ has dimension $n/2-1$. As $x$ is isotropic and orthogonal to it,
$C'=(C_1\cap x^\perp)+\langle x\rangle$ is self-dual, and it is a neighbour of $C_1$. Moreover
$C_1\cap C_2\subseteq C_1\cap x^\perp$, so $C'\cap C_2\supseteq(C_1\cap C_2)+\langle x\rangle$, which is
strictly larger than $C_1\cap C_2$. Induct on $n/2-\dim(C_1\cap C_2)$.
\end{proof}

\begin{theorem}\label{thm:A}
$8\mid\Index(C)$ for every self-dual ternary code $C$.
\end{theorem}

\begin{proof}
By Lemmas~\ref{lem:A-neighbour} and \ref{lem:A-even}, $\Index(C')\equiv-\Index(C)\pmod8$ for
neighbours. The code $T^{n/4}$, a direct sum of tetracodes, has index $0$: every nonzero tetracode
word has weight $3<4$, so no word of the sum is maximal. Lemma~\ref{lem:connected} finishes the
proof.
\end{proof}

With GJF's $12\mid\Index(C)$ (route 0) this proves the Theorem. The same neighbour technique also
shows that the numbers of even and of odd maximal codewords are each divisible by $8$ (project
notes, Proposition~5.5).

\section{Route B: the design identity from modular forms (sketch)}\label{sec:B}

Let $\Lambda=\Lambda(C)=\{x/\sqrt3:\ x\in\Z^n,\ x\bmod3\in C\}$, an odd unimodular lattice
\cite[\S2]{GJF}; $\chi=\sqrt3(1,\dots,1)\in\Lambda$ is characteristic. On the shadow
$S=\Lambda+\chi/2$ put $s(\lambda+\chi/2)=(-1)^{\langle\lambda,\chi\rangle}$. Every $v\in S$ is
$v=y/(2\sqrt3)$ with all $y_k$ odd and $y\bmod3\in C$, so $v^2=|y|^2/12\ge n/12$, with equality iff
$y\in\{\pm1\}^n$. These minimal vectors correspond to the maximal codewords, with
$y=-\varepsilon$ and matching signs.

For $i\ne j$ the polynomial $H(v)=v_iv_j$ is harmonic of degree $2$. Put
\[
  \Phi(\tau)=\sum_{v\in S}s(v)H(v)\,q^{v^2/2}.
\]
Then
\begin{itemize}
\item $\Phi(\tau+1)=e(n/24)\Phi(\tau)$, because $|y|^2\equiv n\pmod{24}$ for $v\in S$;
\item $\Phi(-1/\tau)=(\tau/i)^{n/2}\tau^2\,\Phi(\tau)$, by Poisson summation for theta series with
harmonic coefficients (see e.g.\ \cite{Eb}) and unimodularity, as for the Ramond--Ramond theta series
in \cite[Prop.~2.4]{GJF}.
\end{itemize}
These are the multipliers of $\eta^n\cdot\tau^2$, so $\Phi/\eta^n$ is a modular form of weight $2$
for $\mathrm{SL}_2(\Z)$, holomorphic at $\infty$ because $\Phi=O(q^{n/24})$. Since
$M_2(\mathrm{SL}_2(\Z))=0$ \cite[Ch.~VII]{Se}, $\Phi=0$. Its $q^{n/24}$-coefficient is
$\frac1{12}\sum_{\cF}s(\varepsilon)\varepsilon_i\varepsilon_j$, which proves Theorem~\ref{thm:design}.

\begin{remark}[lattice form]
The argument never uses the code, only $\Lambda$. For every odd unimodular lattice $L$ of rank
$n=12k$ whose shadow has minimum $\ge n/12$, a variant of this proof shows that the signed count of
shadow vectors of norm $n/12$ is divisible by $24$ (project notes, \S2). Route C has no known lattice
analogue.
\end{remark}

\section{Route 0: Gaiotto--Johnson-Freyd for \texorpdfstring{$24\mid n$}{24 | n}}\label{sec:gjf}

For $n=24k$ let $\Lambda^\pm$ be the two even unimodular neighbours of $\Lambda(C)$. Then
$\Index(C)=\Delta^{-k}(\Theta_{\Lambda^+}-\Theta_{\Lambda^-})=a_0^+-a_0^-$, where $a_0^\pm$ is the
$\Delta^k$-coefficient of $\Theta_{\Lambda^\pm}$ in the basis $\Delta^{k-i}c_4^{3i}$. Borcherds
\cite[Thm.~12.1]{Bo} shows $24\mid a_0^\pm$. For $12\mid n$ this gives $24\mid\Index(C\oplus C)
=\Index(C)^2$, hence $12\mid\Index(C)$ \cite[\S1]{GJF}.

\section{Comparison}\label{sec:compare}

\begin{center}
\small
\begin{tabular}{@{}llll@{}}
\toprule
Route & Proves & Main input & Scope\\
\midrule
0 (GJF) & $24\mid\Index$ if $24\mid n$; $12\mid\Index$ & Borcherds' theorem & codes, via lattices\\
A & $8\mid\Index$ & Fourier on $\F^n$, neighbours & codes\\
B & design, hence $8$ and $3$ & harmonic theta series, $M_2=0$ & odd unimodular lattices\\
C & design, hence $8$ and $3$ & one Poisson summation on $\F^n$ & codes\\
C$'$ & design, hence $8$ and $3$ & Poisson, neighbours & codes\\
\bottomrule
\end{tabular}
\end{center}

Route C is the shortest and the only one that is both complete and elementary. Route B is the most
general. Route A gives the extra parity statement mentioned after Theorem~\ref{thm:A}.

\begin{remark}[formal verification]\label{rem:lean}
Route C together with Proposition~\ref{prop:design-suffices} is formalized in Lean~4 with Mathlib
(repository file \texttt{lean/GJF/Divisibility.lean}, about $660$ lines). The formal statement is
\[
  \texttt{index\_dvd\_24}:\quad C=C^\perp\ \text{in}\ \F^n,\ n>0\ \Longrightarrow\ 24\mid
  \#\{\text{even maximal}\}-\#\{\text{odd maximal}\},
\]
with codes as $\F$-submodules of $\F^n$ and the index defined literally as in \cite{GJF}. It uses
only the standard axioms (propositional extensionality, choice, quotient soundness), and all finite
case checks are done by the kernel (\texttt{decide}, no \texttt{native\_decide}). The Fourier
analysis is carried out with $\omega=e^{2\pi i/3}\in\mathbb C$; $|C|^2=3^n$ is obtained by Poisson
summation of the constant function, so no dimension theory is needed. A companion file checks by
\texttt{decide} that the tetracode is self-dual, so the hypotheses are satisfiable.
\end{remark}

\begin{remark}[verification]
Routes A and C, Proposition~\ref{prop:design-suffices} and Remark~\ref{rem:four} were checked by
exact computation (scripts in the project repository); route B only through its conclusion, the
design identity. Route C:
the one-coordinate transforms exactly in $\Z[\omega]$; the pointwise identity of
Lemma~\ref{lem:pointwise} by brute force in $\Z[\omega]$ on every codeword of three codes of length
$12$; $G_{ij}(C)=0$ on random codes of lengths $12$ and $24$; the neighbour form on random neighbour
pairs; for $40$ random non-self-orthogonal codes of length $12$ the correction terms are nonzero
and the conclusion fails in $28$ cases; and Remark~\ref{rem:four} on all $495$ four-sets for two
codes of length $12$ and on $50$ four-sets at length $24$. The design identity and both steps of
Proposition~\ref{prop:design-suffices} hold exactly on random codes of lengths $24$ to $60$, on the
Pless symmetry codes $P(17)$, $P(23)$, $P(29)$ and on the extended quadratic-residue codes
$\mathrm{XQR}(47)$, $\mathrm{XQR}(59)$; for example $|\Index(P(17))|=744$ and
$|\Index(P(29))|=21696$.
\end{remark}

\begin{thebibliography}{9}
\bibitem{GJF}
D.~Gaiotto and T.~Johnson-Freyd,
\emph{Holomorphic SCFTs with small index},
Canad.\ J.\ Math.\ \textbf{74} (2022), no.~2, 573--601. arXiv:1811.00589.

\bibitem{Bo}
R.~E.~Borcherds,
\emph{Automorphic forms on $\mathrm{O}_{s+2,2}(\mathbb R)$ and infinite products},
Invent.\ Math.\ \textbf{120} (1995), 161--213.

\bibitem{Oz}
M.~Ozeki,
\emph{On the notion of Jacobi polynomials for codes},
Math.\ Proc.\ Cambridge Philos.\ Soc.\ \textbf{121} (1997), 15--30.

\bibitem{CM}
H.~S.~Chakraborty and T.~Miezaki,
\emph{Variants of Jacobi polynomials in coding theory},
Des.\ Codes Cryptogr.\ \textbf{90} (2022), 2583--2597.

\bibitem{AHM}
M.~Araya, M.~Harada and K.~Momihara,
\emph{Hadamard matrices related to a certain series of ternary self-dual codes},
arXiv:2205.15498.

\bibitem{Eb}
W.~Ebeling,
\emph{Lattices and Codes}, 3rd ed., Springer Spektrum, 2013.

\bibitem{Se}
J.-P.~Serre,
\emph{A Course in Arithmetic}, Grad.\ Texts in Math.\ \textbf{7}, Springer, 1973.
\end{thebibliography}

\end{document}
